\chapter{ดัชนีคำศัพท์จุลชีววิทยาการเกษตร}
\label{app:glossary}

\begin{description}[font=\bfseries\color{agriGreen}, leftmargin=4.5cm, style=nextline]
    \item[Actinomycetes] แอกทิโนมัยซีท แบคทีเรียแกรมบวกที่มีการเจริญของเส้นใยคล้ายเชื้อราและสร้างสปอร์ เป็นแหล่งผลิตสารปฏิชีวนะสำคัญ
    \item[Ammonification] แอมโมนิฟิเคชัน กระบวนการย่อยสลายสารประกอบไนโตรเจนอินทรีย์ให้กลายเป็นก๊าซแอมโมเนียหรือไอออนแอมโมเนียม
    \item[Antagonism] ปฏิปักษ์ ปรากฏการณ์ที่สิ่งมีชีวิตชนิดหนึ่งยับยั้ง ขัดขวาง หรือทำลายการเจริญเติบโตของสิ่งมีชีวิตอีกชนิดหนึ่ง
    \item[Arbuscule] อาร์บัสคูล โครงสร้างแตกแขนงคล้ายต้นไม้เล็กๆ ของเชื้อราไมคอร์ไรซาภายในเซลล์คอร์เทกซ์รากพืช ใช้แลกเปลี่ยนสารอาหาร
    \item[Biochar] ถ่านชีวภาพ วัสดุคาร์บอนเสถียรที่ได้จากการเผาชีวมวลภายใต้สภาวะควบคุมและจำกัดออกซิเจน
    \item[Biocontrol] การควบคุมโดยชีววิธี การใช้สิ่งมีชีวิตที่เป็นประโยชน์ในการควบคุมและกำจัดศัตรูพืช
    \item[Biological Nitrogen Fixation (BNF)] การตรึงไนโตรเจนทางชีวภาพ การเปลี่ยนก๊าซไนโตรเจนในบรรยากาศเป็นแอมโมเนียโดยเอนไซม์ไนโตรจีเนสของจุลินทรีย์
    \item[Composting] การทำปุ๋ยหมัก กระบวนการย่อยสลายสารอินทรีย์โดยจุลินทรีย์แบบใช้ออกซิเจนจนได้สารปรับปรุงดินที่คงตัว
    \item[Denitrification] ดีไนตริฟิเคชัน กระบวนการรีดักชันของไนเตรตให้กลายเป็นก๊าซไนตรัสออกไซด์หรือก๊าซไนโตรเจนภายใต้สภาวะไร้ออกซิเจน
    \item[Endophyte] เอนโดไฟต์ จุลินทรีย์ที่อาศัยอยู่ภายในเนื้อเยื่อพืชโดยไม่ก่อให้เกิดอันตรายหรือโรคต่อพืชอาศัย
    \item[Glomalin] กลูมาลิน ไกลโคโปรตีนไม่ละลายน้ำที่สร้างโดยเชื้อราไมคอร์ไรซา ทำหน้าที่ยึดเหนี่ยวอนุภาคดินเป็นเม็ดดิน
    \item[Metagenomics] เมตาเจโนมิกส์ การศึกษาลำดับเบสสารพันธุกรรมทั้งหมดของชุมชนจุลินทรีย์ที่สกัดได้โดยตรงจากตัวอย่างธรรมชาติ
    \item[Mycoparasitism] การเป็นปรสิตบนเส้นใยรา ปรากฏการณ์ที่เชื้อราชนิดหนึ่งเข้าเกาะ พันรัด และย่อยสลายเส้นใยของเชื้อราอีกชนิดหนึ่ง
    \item[Necromass] ซากจุลินทรีย์ ซากเซลล์ของจุลินทรีย์ที่ตายแล้วและสะสมเป็นแหล่งคาร์บอนที่เสถียรในดิน
    \item[Nitrification] ไนตริฟิเคชัน กระบวนการออกซิเดชันของแอมโมเนียมให้กลายเป็นไนไตรต์และไนเตรตโดยแบคทีเรียคีโมออโตโทรฟ
    \item[PGPR] แบคทีเรียส่งเสริมการเจริญเติบโตของพืช (Plant Growth-Promoting Rhizobacteria)
    \item[Rhizosphere] ไรโซสเฟียร์ บริเวณดินรอบรากพืชที่ได้รับอิทธิพลโดยตรงจากสารหลั่งและกิจกรรมของราก
    \item[Siderophore] ไซเดอโรฟอร์ สารประกอบน้ำหนักโมเลกุลต่ำที่จุลินทรีย์หลั่งออกไปเพื่อจับและขนส่งไอออนเหล็กเข้าสู่เซลล์
\end{description}

\clearpage